% ============================================================================
% LEHRERHINWEISE: de + artículo (del) + Así se dice
% Klasse 7 | Spanisch | Unidad 2 | DS 03
% ============================================================================

\documentclass[11pt,a4paper]{article}

\usepackage[utf8]{inputenc}
\usepackage[T1]{fontenc}
\usepackage[ngerman]{babel}
\usepackage{geometry}
\usepackage{fancyhdr}
\usepackage{xcolor}
\usepackage{enumitem}
\usepackage{tabularx}
\usepackage{array}
\usepackage{tcolorbox}
\usepackage{booktabs}
\usepackage{amssymb}

% ============================================================================
% LAYOUT
% ============================================================================
\geometry{left=2cm, right=2cm, top=2cm, bottom=2cm}
\definecolor{fach}{HTML}{c41e3a}
\definecolor{grau}{HTML}{666666}
\definecolor{hellgrau}{HTML}{f5f5f5}
\definecolor{basis}{HTML}{2a7a4b}
\definecolor{standard}{HTML}{2c5aa0}
\definecolor{erweiterung}{HTML}{7b1fa2}
\definecolor{spiel}{HTML}{b35c00}

\pagestyle{fancy}
\fancyhf{}
\fancyhead[L]{\small\textcolor{grau}{Spanisch Klasse 7 -- Unidad 2}}
\fancyhead[R]{\small\textcolor{grau}{LH-03 (Lehrerhinweise)}}
\fancyfoot[C]{\small\thepage}
\renewcommand{\headrulewidth}{0.4pt}

% Spanisch-Befehl
\newcommand{\es}[1]{\textbf{#1}}

% ============================================================================
% DOKUMENT
% ============================================================================
\begin{document}

% ----------------------------------------------------------------------------
% KOPFBEREICH
% ----------------------------------------------------------------------------
\begin{center}
{\Large\textcolor{fach}{\textbf{Lehrerhinweise}}}\\[0.3cm]
{\large de + artículo (del) + Así se dice}\\[0.2cm]
{\normalsize Doppelstunde (90 Min.) | Qué pasa 1, S. 42--43}
\end{center}

\vspace{0.5cm}
\hrule
\vspace{0.5cm}

% ----------------------------------------------------------------------------
% KURZÜBERSICHT
% ----------------------------------------------------------------------------
\section*{Kurzübersicht}

\begin{tabular}{ll}
\textbf{Thema:} & Kontraktion de+el=del + Klassenraumredemittel \\
\textbf{Lehrwerk:} & Qué pasa 1, Unidad 2, S. 42--43 \\
\textbf{Materialien:} & AB-03, LB-03, ML-03, Backup-Box, Ampelkarten \\
\textbf{Fokus:} & Besitz ausdrücken (de+artículo), Así se dice \\
\textbf{Voraussetzung:} & hay + estar (DS 02), bestimmte Artikel \\
\end{tabular}

\vspace{0.5cm}

% ----------------------------------------------------------------------------
% FEINZIELE
% ----------------------------------------------------------------------------
\section*{Feinziele}

\begin{enumerate}[label=\textbf{FZ\arabic*:}, leftmargin=2.5em]
    \item SuS \textbf{wiederholen} hay vs estar in Kontextsätzen \textit{(AFB I)} -- Kurz-Repaso
    \item SuS \textbf{bilden} die Kontraktion de + el = del korrekt \textit{(AFB I)}
    \item SuS \textbf{drücken} Besitz mit de + artículo aus: \es{la silla de la profesora} \textit{(AFB II)}
    \item SuS \textbf{verwenden} Klassenraumredemittel: \es{¿Cómo?}, \es{No comprendo}, \es{Más despacio...} \textit{(AFB I)}
\end{enumerate}

\vspace{0.5cm}

% ----------------------------------------------------------------------------
% STUNDENVERLAUF
% ----------------------------------------------------------------------------
\section*{Stundenverlauf (90 Min.)}

\begin{tabular}{|p{1cm}|p{2.2cm}|p{6cm}|p{1cm}|p{2.8cm}|}
\hline
\textbf{Min} & \textbf{Phase} & \textbf{Aktivität} & \textbf{SF} & \textbf{Material/Signal} \\
\hline
0--5 & Einstieg & Repaso hay/estar: 4 Sätze, Daumen hoch/runter & PL & Tafel \\
\hline
5--15 & Erarbeitung 1 & Input de+artículo: 3 Beispiele an Tafel, Regel entdecken & PL & Tafel, LB-03, \newline\textbf{Signal:} Gong \\
\hline
15--25 & Übung 1 & AB Aufg. 1: Zuordnung, dann Aufg. 2: Lücken & EA & AB-03, Timer 10' \\
\hline
25--40 & \textbf{Spiel} & \es{¿De quién es?} mit Backup-Box & GA/PL & Box, Tafel-Hilfe, \newline\textbf{Signal:} Musik aus \\
\hline
40--45 & Pause & Bewegung: Aufstehen, strecken & -- & -- \\
\hline
45--55 & Erarbeitung 2 & Así se dice: Redemittel einführen, chorisch 3× & PL & LB-03, Poster \\
\hline
55--65 & Drill & Redemittel-Drill mit Ampelkarten & PL/PA & Karten, \newline\textbf{Signal:} Grün \\
\hline
65--78 & Anwendung & AB Aufg. 3--4: hay/estar + Minidialog & PA & AB-03, Timer 13' \\
\hline
78--85 & Sicherung & Merkregel: \es{de + el = del, sonst getrennt} & PL & Tafel \\
\hline
85--90 & Puffer & Redemittel-Poster fertig / Aufräumen & PL & Poster, Gong \\
\hline
\end{tabular}

\vspace{0.5cm}

% ----------------------------------------------------------------------------
% DIFFERENZIERUNG
% ----------------------------------------------------------------------------
\section*{Differenzierung}

\begin{tcolorbox}[colback=white, colframe=basis, title=\textcolor{white}{\textbf{[B] Basis}}]
\begin{itemize}[nosep]
    \item Aufg. 1: Zuordnung MIT farbiger Hilfe (del = rot markiert)
    \item Aufg. 2: Lückentext mit Wortbank (del, de la, de los, de las)
    \item Aufg. 3: hay/estar nur 2 Sätze
    \item Spiel: Darf Tafel-Hilfe benutzen, wird nicht als erstes drangenommen
    \item Redemittel: Karte mit allen Phrasen vor sich
\end{itemize}
\end{tcolorbox}

\begin{tcolorbox}[colback=white, colframe=standard, title=\textcolor{white}{\textbf{[S] Standard}}]
\begin{itemize}[nosep]
    \item Aufg. 1--4: Vollständig ohne Hilfe
    \item Spiel: Normale Teilnahme
    \item Redemittel: Nur Poster als Hilfe
\end{itemize}
\end{tcolorbox}

\begin{tcolorbox}[colback=white, colframe=erweiterung, title=\textcolor{white}{\textbf{[E] Erweiterung}}]
\begin{itemize}[nosep]
    \item Aufg. 1--4: Schnell fertig $\rightarrow$ Zusatz: 4 eigene Besitz-Sätze
    \item Spiel: Übernimmt Spielleitung (zeigt Gegenstände)
    \item Redemittel: Muss auch auf unbekannte Situation reagieren (L flüstert)
\end{itemize}
\end{tcolorbox}

\vspace{0.5cm}

% ----------------------------------------------------------------------------
% SPIELANLEITUNG
% ----------------------------------------------------------------------------
\section*{Spiel: ¿De quién es?}

\begin{tcolorbox}[colback=white, colframe=spiel, title=\textcolor{white}{\textbf{Spielablauf (15 Min.)}}]

\textbf{Material:} Backup-Box mit 15 Gegenständen (falls SuS vergessen)

\vspace{0.3cm}

\textbf{Vorbereitung (2 Min):}
\begin{enumerate}[nosep]
    \item L stellt Backup-Box auf Lehrertisch
    \item SuS legen 1 Gegenstand dazu (wer keinen hat $\rightarrow$ Box)
    \item Alle Gegenstände durchmischen
\end{enumerate}

\vspace{0.3cm}

\textbf{Ablauf:}
\begin{enumerate}[nosep]
    \item L nimmt Gegenstand: \es{¿De quién es este bolígrafo?}
    \item SuS melden sich, L nimmt dran
    \item SuS rät: \es{Es el bolígrafo de María.}
    \item Besitzer: \es{¡Sí!} oder \es{No, no es mi bolígrafo.}
    \item Bei richtig: Ratender darf nächsten Gegenstand nehmen
    \item \textbf{Tempo-Regel:} Max. 30 Sek pro Gegenstand
\end{enumerate}

\vspace{0.3cm}

\textbf{Sprachgerüst (an Tafel):}
\begin{itemize}[nosep]
    \item \es{Es el/la ... de [Name].}
    \item \es{Es el/la ... del profesor.}
    \item \es{Es el/la ... de la profesora.}
    \item \es{¡Sí, es mi ...!} / \es{No, no es mi ...}
\end{itemize}

\vspace{0.3cm}

\textbf{Notfallplan (falls Spiel nicht funktioniert):}

$\rightarrow$ Umwandeln in \es{Veo, veo...} mit del/de la:
\es{Veo algo del profesor...} -- SuS raten was.

\end{tcolorbox}

\vspace{0.5cm}

% ----------------------------------------------------------------------------
% REDEMITTEL-DRILL
% ----------------------------------------------------------------------------
\section*{Redemittel-Drill mit Ampelkarten}

\begin{tcolorbox}[colback=white, colframe=fach, title=\textcolor{white}{\textbf{Así se dice -- Drill (10 Min.)}}]

\textbf{Material:} Ampelkarten (pro SuS: 1× grün, 1× rot)

\vspace{0.3cm}

\textbf{Funktion:}
\begin{itemize}[nosep]
    \item \textbf{Grün} = Jetzt reagieren!
    \item \textbf{Rot} = Nur zuhören
\end{itemize}

\vspace{0.3cm}

\textbf{Ablauf:}
\begin{enumerate}[nosep]
    \item L erklärt: \glqq Wenn ich GRÜN zeige, müsst ihr mit einem Redemittel reagieren!\grqq
    \item L spricht normal $\rightarrow$ ROT $\rightarrow$ SuS hören zu
    \item L spricht zu schnell $\rightarrow$ GRÜN $\rightarrow$ SuS: \es{¡Más despacio, por favor!}
    \item L sagt unbekanntes Wort $\rightarrow$ GRÜN $\rightarrow$ SuS: \es{¿Qué significa...?}
    \item L flüstert $\rightarrow$ GRÜN $\rightarrow$ SuS: \es{¿Cómo?}
\end{enumerate}

\vspace{0.3cm}

\textbf{Steigerung:}
\begin{itemize}[nosep]
    \item Runde 1: L zeigt Karte (leicht)
    \item Runde 2: L zeigt keine Karte, SuS müssen selbst erkennen (schwer)
    \item Runde 3: Partner-Drill (A flüstert, B reagiert)
\end{itemize}

\end{tcolorbox}

\vspace{0.5cm}

% ----------------------------------------------------------------------------
% ÜBERGANGSSIGNALE
% ----------------------------------------------------------------------------
\section*{Übergangssignale}

\begin{tabular}{|l|l|l|}
\hline
\textbf{Von $\rightarrow$ Nach} & \textbf{Signal} & \textbf{Ansage} \\
\hline
Einstieg $\rightarrow$ Erarbeitung & Gong & \es{Libros a la página 42} \\
\hline
Erarbeitung $\rightarrow$ Übung & Timer Start & \es{Tenéis 10 minutos} \\
\hline
Übung $\rightarrow$ Spiel & Musik & \es{¡Levantaos! Juego.} \\
\hline
Spiel $\rightarrow$ Pause & Musik aus & \es{Pausa, 5 minutos} \\
\hline
Pause $\rightarrow$ Erarbeitung 2 & Gong & \es{Sentaos, miramos el póster} \\
\hline
Drill $\rightarrow$ Anwendung & Grüne Karte hoch & \es{Ahora con pareja} \\
\hline
Anwendung $\rightarrow$ Sicherung & Timer Ende & \es{Stop, miramos aquí} \\
\hline
\end{tabular}

\vspace{0.5cm}

% ----------------------------------------------------------------------------
% MATERIALCHECKLISTE
% ----------------------------------------------------------------------------
\section*{Materialcheckliste}

\textbf{Vom Lehrer vorzubereiten:}
\begin{itemize}[nosep]
    \item[$\square$] \textbf{Backup-Box} mit 15 Gegenständen:
    \begin{itemize}[nosep]
        \item 3× Stifte (verschiedene Farben)
        \item 2× Radiergummi, 2× Lineal, 2× kleines Heft
        \item 2× Schere, 2× Anspitzer, 2× Klebestift
    \end{itemize}
    \item[$\square$] \textbf{Ampelkarten} (pro SuS): 1× grün, 1× rot (Pappe, 5×5cm)
    \item[$\square$] \textbf{Redemittel-Poster} (A2, vorgedruckt oder Tafel vorbereitet)
    \item[$\square$] \textbf{Timer} (Handy oder Küchenuhr, sichtbar)
    \item[$\square$] \textbf{Gong/Klingel} für Phasenwechsel
\end{itemize}

\vspace{0.3cm}

\textbf{Kopien:}
\begin{itemize}[nosep]
    \item[$\square$] AB-03 (Klassensatz)
    \item[$\square$] LB-03 (Klassensatz)
    \item[$\square$] Lehrwerk Qué pasa 1 (S. 42--43)
\end{itemize}

\vspace{0.3cm}

\textbf{Von SuS mitzubringen:}
\begin{itemize}[nosep]
    \item 1--2 Schulgegenstände (Backup vorhanden!)
    \item AB aus letzter Stunde (hay/estar)
\end{itemize}

\vspace{0.5cm}

% ----------------------------------------------------------------------------
% HÄUFIGE FEHLER
% ----------------------------------------------------------------------------
\section*{Häufige Fehler und Reaktionen}

\begin{tabular}{|p{3.5cm}|p{3.5cm}|p{6cm}|}
\hline
\textbf{Fehler} & \textbf{Ursache} & \textbf{Reaktion} \\
\hline
*\es{de el profesor} & Kontraktion vergessen & Immer: de + el = \textbf{del}! \\
\hline
*\es{del profesora} & Übergeneralisierung & del nur maskulin! \es{de la} profesora \\
\hline
*\es{del los alumnos} & Plural-Verwechslung & Nur Singular: del. Plural: \es{de los} \\
\hline
*\es{la libro de María} & Genus-Fehler & \es{el libro} -- maskulin! \\
\hline
\es{Como?} ohne ¿ & Satzzeichen vergessen & Im Spanischen: ¿...? \\
\hline
\end{tabular}

\vspace{0.5cm}

% ----------------------------------------------------------------------------
% TAFELBILD
% ----------------------------------------------------------------------------
\section*{Tafelbild}

\begin{tcolorbox}[colback=white, colframe=grau]
\textbf{de + artículo}

\vspace{0.3cm}
\begin{tabular}{p{6cm}|p{6cm}}
\textbf{Kontraktion (nur 1 Fall!)} & \textbf{Keine Kontraktion} \\
\hline
de + el = \textbf{del} & de + la = de la \\
\es{el libro \textbf{del} alumno} & \es{la silla \textbf{de la} profesora} \\
& de + los = de los \\
& de + las = de las \\
\end{tabular}

\vspace{0.5cm}

\textbf{Así se dice:}

\vspace{0.2cm}
\begin{tabular}{ll}
\es{¿Cómo?} & Wie bitte? \\
\es{No comprendo.} & Ich verstehe nicht. \\
\es{¿Qué significa...?} & Was bedeutet...? \\
\es{Más despacio, por favor.} & Langsamer, bitte. \\
\es{¿Puedes repetir?} & Kannst du wiederholen? \\
\end{tabular}
\end{tcolorbox}

\vspace{0.5cm}

% ----------------------------------------------------------------------------
% AB-LÖSUNGEN
% ----------------------------------------------------------------------------
\section*{AB-Lösungen (Kurzform)}

\textbf{Aufg. 1 -- Zuordnung:}
\begin{itemize}[nosep]
    \item de + el $\rightarrow$ \textbf{del}
    \item de + la $\rightarrow$ de la
    \item de + los $\rightarrow$ de los
    \item de + las $\rightarrow$ de las
\end{itemize}

\vspace{0.2cm}

\textbf{Aufg. 2 -- Lücken:}
a) del alumno, b) de la profesora, c) de los alumnos, d) de las alumnas, e) del profesor, f) de la clase

\vspace{0.2cm}

\textbf{Aufg. 3 -- hay/estar:}
a) hay, b) está, c) hay, d) están, e) hay, f) está

\vspace{0.2cm}

\textbf{Aufg. 4 -- Minidialog:} Individuelle Lösungen mit Redemitteln

\vspace{0.5cm}

\end{document}
